\subsubsection{Binary Lifting / LCA}
\begin{lstlisting}[style=mystyle, language=C++]
// TC: O(N log N) build, O(log N) per LCA query
// usa: lowest common ancestor en arbol estatico (rooted)
#include <bits/stdc++.h>
using namespace std;

struct LCA{
	int n, LOG;
	vector<vector<int>> up;
	vector<int> depth;
	vector<vector<int>> adj;
	LCA(int n): n(n){
		LOG = 32 - __builtin_clz(n + 1);
		up.assign(LOG, vector<int>(n, -1));
		depth.assign(n, 0);
		adj.assign(n, {});
	}
	void addEdge(int u, int v){
		adj[u].push_back(v);
		adj[v].push_back(u);
	}
	void dfs(int v, int p){
		up[0][v] = p;
		for(int k=1;k<LOG;k++){
			if(up[k-1][v] != -1) 
				up[k][v] = up[k-1][up[k-1][v]];
		}
		for(int u : adj[v]){
			if(u == p) continue;
			depth[u] = depth[v] + 1;
			dfs(u, v);
		}
	}
	void build(int root){
		dfs(root, -1);
	}
	int lca(int a, int b){
		if(depth[a] < depth[b]) swap(a, b);
		int diff = depth[a] - depth[b];
		for(int k=0;k<LOG;k++) if(diff & (1<<k)) a = up[k][a];
		if(a == b) return a;
		for(int k=LOG-1;k>=0;k--){
			if(up[k][a] != up[k][b]){
				a = up[k][a];
				b = up[k][b];
			}
		}
		return up[0][a];
	}
};

int main(){
	LCA lca(5);
	lca.addEdge(0,1); lca.addEdge(1,2); lca.addEdge(1,3); lca.addEdge(3,4);
	lca.build(0);
	cout << "LCA(2,4) = " << lca.lca(2,4) << "\n";  // 1
	return 0;
}
\end{lstlisting}
